\documentclass[11pt,a4paper]{article}
\usepackage{amsmath,amssymb,amsthm}
\usepackage[margin=28mm]{geometry}
\usepackage[colorlinks=true,linkcolor=blue,urlcolor=blue]{hyperref}
\input{glyphtounicode}
\pdfgentounicode=1
\hypersetup{
  pdftitle={Uniform geometric estimates and Mandelbrot local connectivity},
  pdfsubject={Conditional MLC theorem from an unproved uniform geometric hypothesis}
}

\newtheorem{theorem}{Theorem}[section]
\newtheorem{lemma}[theorem]{Lemma}
\newtheorem{proposition}[theorem]{Proposition}
\theoremstyle{definition}
\newtheorem{definition}[theorem]{Definition}

\newcommand{\C}{\mathbb C}
\newcommand{\R}{\mathbb R}
\newcommand{\N}{\mathbb N}
\newcommand{\M}{\mathcal M}
\newcommand{\D}{D}
\newcommand{\id}{\operatorname{id}}
\newcommand{\LC}{\operatorname{LC}}
\newcommand{\Conn}{\operatorname{Connected}}

\title{Uniform geometric estimates\\and Mandelbrot local connectivity}
\author{}
\date{}

\begin{document}
\maketitle

\section*{Task}
Prove $\LC(\M)$ by establishing the existence assertion $\mathsf{UG}$
in Definition~\ref{def:ug}.  Theorem~\ref{thm:mlc} proves
$\mathsf{UG}\Longrightarrow\LC(\M)$.  The assertion $\mathsf{UG}$
remains unproved.

\section{Definitions and local connectedness}

All subsets of $\C$ carry the Euclidean subspace topology.
The notation $\Conn(V)$ asserts that $V$ is connected.

\begin{definition}[Critical orbit and Mandelbrot set]\label{def:mandelbrot}
\begin{gather*}
 \N:=\{0,1,2,\ldots\},\qquad f_c(z):=z^2+c
 \quad(c,z\in\C),\\
 p_0(c):=0,\qquad p_{n+1}(c):=p_n(c)^2+c
 \quad(c\in\C,\ n\in\N),\\
 \M:=\{c\in\C:\exists B\in[0,\infty)\ \forall n\in\N,\
                         |p_n(c)|\le B\}.
\end{gather*}
\end{definition}
\noindent
In words, $\N$ is the set of nonnegative integers.  For each complex
parameter $c$, $f_c$ sends $z$ to $z^2+c$.  The critical orbit starts
at $p_0(c)=0$ and satisfies $p_{n+1}(c)=p_n(c)^2+c$ for every
$n\in\N$.  A complex parameter $c$ belongs to $\M$ if and only if
there is a nonnegative real number $B$ such that $|p_n(c)|\le B$
for every $n\in\N$.

\begin{definition}[Connected neighborhoods]\label{def:lc}
For $c\in\C$, $r>0$, and $X\subseteq\C$, define
\begin{gather*}
 B(c,r):=\{z\in\C:|z-c|<r\},\qquad
 \overline B(c,r):=\{z\in\C:|z-c|\le r\},\\
 \LC(X):=
 \left[
 \begin{gathered}
 \forall c\in X\ \forall\eta>0\ \exists\delta>0\
 \exists V\subseteq X:\\
 \Conn(V)\ \land\
 X\cap B(c,\delta)\subseteq V\subseteq X\cap B(c,\eta)
 \end{gathered}
 \right].
\end{gather*}
\end{definition}
\noindent
In words, $B(c,r)$ consists of the points whose distance from $c$ is
less than $r$, and $\overline B(c,r)$ consists of the points whose
distance from $c$ is at most $r$.  The assertion $\LC(X)$ holds if and
only if, for every $c\in X$ and every positive real number $\eta$,
there exist a positive real number $\delta$ and a connected subset
$V$ of $X$ containing $X\cap B(c,\delta)$ and contained in
$X\cap B(c,\eta)$.

\begin{lemma}[Connected-neighborhood criterion]\label{lem:lc-criterion}
For every $X\subseteq\C$,
\[
 \LC(X)\quad\Longleftrightarrow\quad
 \text{$X$ has a basis of connected relatively open sets}.
\]
\end{lemma}
\noindent
In words, $\LC(X)$ holds if and only if, for every relatively open
subset $U$ of $X$ and every $c\in U$, there is a connected relatively
open subset $V$ of $X$ with $c\in V\subseteq U$.
\begin{proof}
Suppose that $X$ has a basis of connected relatively open sets.
For $c\in X$ and $\eta>0$, choose a connected relatively open set $V$
such that
\[
 c\in V\subseteq X\cap B(c,\eta).
\]
Relative openness gives $\delta>0$ with
$X\cap B(c,\delta)\subseteq V$.  These choices satisfy
Definition~\ref{def:lc}.

Conversely, assume $\LC(X)$.  Fix a relatively open subset $U$ of
$X$ and $c\in U$, and define
\[
 K:=\bigcup\{S\subseteq U:\Conn(S)\ \land\ c\in S\}.
\]
The family contains $\{c\}$.  Every member is connected and contains
$c$, so $K$ is connected and $c\in K\subseteq U$.
For any $y\in K$, relative openness of $U$ gives $\eta>0$ such that
$X\cap B(y,\eta)\subseteq U$.  Applying $\LC(X)$ at $y$ gives
$\delta>0$ and a connected set $V_y$ with
\[
 X\cap B(y,\delta)\subseteq V_y
 \subseteq X\cap B(y,\eta)\subseteq U.
\]
The sets $K$ and $V_y$ both contain $y$.  Their union is therefore a
connected subset of $U$ containing $c$, so the definition of $K$ gives
$K\cup V_y\subseteq K$.  In particular,
$X\cap B(y,\delta)\subseteq K$.  This holds for every $y\in K$;
hence $K$ is relatively open in $X$.  The set $K$ is the required
connected relatively open neighborhood of $c$ inside $U$.
\end{proof}

\section{Orbit bounds and outer sets}

\begin{lemma}[Escape radius]\label{lem:escape}
\begin{equation}\label{eq:escape}
 \M=\{c\in\C:\forall n\in\N,\ |p_n(c)|\le2\}.
\end{equation}
\end{lemma}
\noindent
In words, a complex parameter $c$ belongs to $\M$ if and only if
$|p_n(c)|\le2$ for every $n\in\N$.
\begin{proof}
First suppose that $c\in\C$ and $n\in\N$ satisfy
\[
 q:=|p_n(c)|>2,\qquad |c|\le q.
\]
Induction on $j$ gives
\begin{equation}\label{eq:escape-growth}
 |p_{n+j}(c)|\ge q(q-1)^j\qquad(j\in\N).
\end{equation}
The case $j=0$ is equality.  For the induction step,
$|p_{n+j}(c)|\ge q$ because $q-1>1$, and consequently
\begin{align*}
 |p_{n+j+1}(c)|
 &\ge |p_{n+j}(c)|^2-|c|\\
 &\ge |p_{n+j}(c)|^2-|p_{n+j}(c)|\\
 &\ge(q-1)|p_{n+j}(c)|\\
 &\ge q(q-1)^{j+1}.
\end{align*}
Since $q-1>1$, the lower bound in \eqref{eq:escape-growth} tends to
infinity.  Thus $c\notin\M$.

If $|c|>2$, apply this argument with $n=1$ and $q=|p_1(c)|=|c|$.
If $|c|\le2$ and $|p_n(c)|>2$ for some $n$, apply it with
$q=|p_n(c)|$, which satisfies $|c|\le q$.  Hence membership in $\M$
implies $|p_n(c)|\le2$ for every $n$.  Conversely, these inequalities
give the bound $B=2$ in Definition~\ref{def:mandelbrot}.
\end{proof}

\begin{definition}[Outer sets and parameter square]\label{def:outer}
For $N\in\N$, define
\begin{equation}\label{eq:outer}
 O_N:=\overline B(0,2)\cap
       \bigcap_{j=0}^{N}p_j^{-1}\bigl(\overline B(0,2)\bigr),
 \qquad
 \D:=\{x+iy:x,y\in[-2,2]\}.
\end{equation}
\end{definition}
\noindent
In words, $O_N$ consists of the parameters $c$ satisfying $|c|\le2$
and $|p_j(c)|\le2$ for every integer $j$ with $0\le j\le N$.
The set $\D$ consists of the complex numbers whose real and imaginary
parts both belong to $[-2,2]$.

\begin{lemma}[Nested closed outer sets]\label{lem:outer}
For every $N\in\N$,
\begin{equation}\label{eq:outer-properties}
 0\in\M\subseteq O_{N+1}\subseteq O_N\subseteq\D,
 \qquad
 \bigcap_{N\in\N}O_N=\M.
\end{equation}
Each $O_N$, $\M$, and $\D$ is compact and closed in $\C$, and $\D$
is convex.
\end{lemma}
\noindent
In words, $0$ belongs to $\M$; every outer set contains $\M$;
$O_{N+1}$ is contained in $O_N$; and every $O_N$ is contained in
$\D$.  The intersection of all outer sets is $\M$.  Each outer set,
the Mandelbrot set, and the parameter square is compact and closed,
and the parameter square is convex.
\begin{proof}
Induction using the orbit recurrence proves both that $p_n$ is
continuous on $\C$ and that $p_n(0)=0$ for every $n\in\N$.
Thus $0\in\M$.  Each $O_N$ is a finite intersection of closed sets
and is contained in $\overline B(0,2)$; it is compact by the
Heine--Borel theorem.  The identity
\[
 O_{N+1}=O_N\cap p_{N+1}^{-1}\bigl(\overline B(0,2)\bigr)
\]
gives nesting.  Lemma~\ref{lem:escape}, together with $p_1(c)=c$,
gives $\M\subseteq O_N$ for every $N$ and
\[
 \bigcap_{N\in\N}O_N
 =\overline B(0,2)\cap
   \bigcap_{j\in\N}p_j^{-1}\bigl(\overline B(0,2)\bigr)
 =\M.
\]
This intersection is closed and bounded, so $\M$ is compact.
The inequalities $|\operatorname{Re}z|\le|z|$ and
$|\operatorname{Im}z|\le|z|$ give $\overline B(0,2)\subseteq\D$.
The square $\D$ is closed and bounded, hence compact.  If $z,w\in\D$
and $t\in[0,1]$, each coordinate of $(1-t)z+tw$ lies in $[-2,2]$.
Thus $\D$ is convex.
\end{proof}

\begin{definition}[Hausdorff outer gap]\label{def:outer-gap}
For nonempty compact sets $A,B\subseteq\C$, define
\begin{gather*}
 \operatorname{dist}(z,A):=\inf_{a\in A}|z-a|,\\
 d_H(A,B):=
 \max\left\{\sup_{a\in A}\operatorname{dist}(a,B),\
             \sup_{b\in B}\operatorname{dist}(b,A)\right\},\\
 h_N:=d_H(O_N,\M)
     =\sup_{z\in O_N}\operatorname{dist}(z,\M)
 \qquad(N\in\N).
\end{gather*}
\end{definition}
\noindent
In words, $h_N$ is the largest distance from a point of the finite-stage
outer set $O_N$ to the Mandelbrot set $\M$.  The second expression for
$h_N$ is valid because $\M\subseteq O_N$.

\begin{theorem}[Hausdorff outer-gap convergence]\label{thm:outer-gap}
\begin{gather}
 0\le h_{N+1}\le h_N,\qquad
 \lim_{N\to\infty}h_N=0,\label{eq:outer-gap-limit}\\
 \exists\nu\in\N^\N:\quad
 \nu_0=0,\qquad
 \forall n\ge1,\quad
 \nu_n>\nu_{n-1},\quad
 \nu_n\ge n,\quad
 h_{\nu_n}\le2^{-n}.\label{eq:depth-selection}
\end{gather}
\end{theorem}
\noindent
In words, the finite-stage outer sets decrease to $\M$ in Hausdorff
distance.  In particular, one can choose a strictly increasing sequence
of orbit depths $\nu_n$ with $\nu_n\ge n$ and with outer-set error at
most $2^{-n}$ for every $n\ge1$.
\begin{proof}
Since $O_{N+1}\subseteq O_N$,
\[
 h_{N+1}
 =\sup_{z\in O_{N+1}}\operatorname{dist}(z,\M)
 \le\sup_{z\in O_N}\operatorname{dist}(z,\M)
 =h_N.
\]
Thus $(h_N)_{N\in\N}$ has a limit $h_\infty\ge0$.

Assume for contradiction that $h_\infty>0$.  For every $N\in\N$,
compactness of $O_N$ and continuity of
$z\mapsto\operatorname{dist}(z,\M)$ give
$z_N\in O_N$ with
\[
 \operatorname{dist}(z_N,\M)=h_N.
\]
The sequence $(z_N)_{N\in\N}$ lies in the compact set $O_0$.
Choose a convergent subsequence
$z_{N_k}\to z_\infty\in O_0$.  Fix $m\in\N$.  For all sufficiently
large $k$, $N_k\ge m$, and hence
\[
 z_{N_k}\in O_{N_k}\subseteq O_m.
\]
The set $O_m$ is closed, so $z_\infty\in O_m$.  Since $m$ was arbitrary,
\[
 z_\infty\in\bigcap_{m\in\N}O_m=\M.
\]
The distance function to a nonempty set is $1$-Lipschitz:
\[
 \left|\operatorname{dist}(z,A)-\operatorname{dist}(w,A)\right|
 \le |z-w|.
\]
Therefore
\[
 h_{N_k}
 =\operatorname{dist}(z_{N_k},\M)
 \longrightarrow\operatorname{dist}(z_\infty,\M)=0,
\]
contradicting $h_{N_k}\to h_\infty>0$.  Hence
$h_N\to0$.

Choose $\nu_0:=0$.  Inductively, suppose that $\nu_{n-1}$ has been
chosen for some $n\ge1$.  Since $h_N\to0$, there exists
$N>\max\{n,\nu_{n-1}\}$ with $h_N\le2^{-n}$.  Set $\nu_n:=N$.
This gives \eqref{eq:depth-selection}.
\end{proof}

\begin{proposition}[Image error from the outer gap]\label{prop:image-gap}
Let $\nu\in\N^\N$ and let
$T_\bullet\in C(\D,\D)^\N$.  Assume
\[
 \forall n\in\N\ \forall c\in O_{\nu_n},\qquad T_n(c)=c.
\]
Define
\[
 a_n:=\sup_{z\in\D}
       \operatorname{dist}(T_n(z),O_{\nu_n}).
\]
Then
\begin{equation}\label{eq:image-gap}
 d_H(T_n(\D),\M)\le a_n+h_{\nu_n}
 \qquad(n\in\N).
\end{equation}
\end{proposition}
\noindent
In words, the Hausdorff distance from the image $T_n(\D)$ to $\M$ is
bounded by the sum of the distance from $T_n(\D)$ to its finite-stage
outer set and the distance from that outer set to $\M$.
\begin{proof}
Fix $n\in\N$.  Since $\M\subseteq O_{\nu_n}$ and $T_n$ fixes
$O_{\nu_n}$,
\[
 \M\subseteq T_n(\M)\subseteq T_n(\D).
\]
Thus
\[
 \sup_{c\in\M}\operatorname{dist}(c,T_n(\D))=0.
\]
Fix $z\in\D$ and $\varepsilon>0$.  Choose
$w\in O_{\nu_n}$ with
\[
 |T_n(z)-w|
 \le\operatorname{dist}(T_n(z),O_{\nu_n})+\varepsilon
 \le a_n+\varepsilon.
\]
Then
\[
 \operatorname{dist}(T_n(z),\M)
 \le |T_n(z)-w|+\operatorname{dist}(w,\M)
 \le a_n+\varepsilon+h_{\nu_n}.
\]
Letting $\varepsilon\downarrow0$ and taking the supremum over
$z\in\D$ gives
\[
 \sup_{y\in T_n(\D)}\operatorname{dist}(y,\M)
 \le a_n+h_{\nu_n}.
\]
Together with the first directed distance, this is
\eqref{eq:image-gap}.
\end{proof}

\section{Uniform geometric hypothesis}

\begin{definition}[Uniform distance]\label{def:uniform-distance}
For topological spaces $X,Y$, define
\begin{gather*}
 C(X,Y):=\{F:X\to Y:F\text{ is continuous}\},\\
 \|F-G\|_\infty:=\sup_{z\in\D}|F(z)-G(z)|
 \qquad(F,G\in C(\D,\C)).
\end{gather*}
\end{definition}
\noindent
In words, $C(X,Y)$ is the set of continuous maps from $X$ to $Y$,
and $\|F-G\|_\infty$ is the supremum of the distances between $F(z)$
and $G(z)$ as $z$ ranges over $\D$.  This supremum is finite because
$\D$ is nonempty and compact and $z\mapsto|F(z)-G(z)|$ is continuous.

\begin{definition}[Uniform geometric existence]\label{def:ug}
For $\nu\in\N^\N$ and $R_\bullet\in C(\D,\D)^\N$, define
\begin{equation}\label{eq:hypotheses}
\begin{aligned}
 \mathsf A(\nu,R_\bullet):={}&(R_0=\id_\D)\\
 &{}\land(\forall n\in\N,\ \nu_n\ge n)\\
 &{}\land(\forall n\in\N\ \forall c\in O_{\nu_n},\ R_n(c)=c)\\
 &{}\land\bigl(\forall n\in\N\ \forall z\in\D\
        \exists w\in O_{\nu_n},\ |R_n(z)-w|\le2^{-n}\bigr)\\
 &{}\land\bigl(\forall n\in\N,\
        \|R_{n+1}-R_n\|_\infty\le8\,2^{-n}\bigr).
\end{aligned}
\end{equation}
\begin{equation}\label{eq:ug-existence}
 \mathsf{UG}:=
 \exists\nu\in\N^\N\ \exists R_\bullet\in C(\D,\D)^\N,\
 \mathsf A(\nu,R_\bullet).
\end{equation}
\end{definition}
\noindent
In words, $\mathsf A(\nu,R_\bullet)$ holds if and only if all the
following conditions hold.  The map $R_0$ is the identity on $\D$.
For every $n\in\N$, the integer $\nu_n$ is at least $n$.
For every $n\in\N$ and every $c\in O_{\nu_n}$, $R_n(c)=c$.
For every $n\in\N$ and every $z\in\D$, there is
$w\in O_{\nu_n}$ with $|R_n(z)-w|\le2^{-n}$.
For every $n\in\N$, the uniform distance between $R_{n+1}$ and $R_n$
is at most $8\,2^{-n}$.

\noindent
In words, $\mathsf{UG}$ holds if and only if there exist a sequence
of natural numbers $\nu$ and a sequence of continuous maps
$R_n:\D\to\D$ for which $\mathsf A(\nu,R_\bullet)$ holds.

\begin{theorem}[Summable candidate maps imply $\mathsf{UG}$]
\label{thm:candidate-to-ug}
Let $\nu\in\N^\N$ and $T_\bullet\in C(\D,\D)^\N$ satisfy
\begin{gather}
 \nu_0=0,\qquad
 \forall n\ge1,\quad
 \nu_n>\nu_{n-1},\quad
 \nu_n\ge n,\quad
 h_{\nu_n}\le2^{-n},\label{eq:candidate-depths}\\
 \forall n\in\N\ \forall c\in O_{\nu_n},\qquad T_n(c)=c,\label{eq:candidate-fixing}\\
 a_n:=\sup_{z\in\D}
       \operatorname{dist}(T_n(z),O_{\nu_n})
 \longrightarrow0,\label{eq:candidate-error}\\
 b_n:=\|T_{n+1}-T_n\|_\infty,\qquad
 \sum_{n=0}^{\infty}b_n<\infty.\label{eq:candidate-summability}
\end{gather}
Then $\mathsf{UG}$ holds.
\end{theorem}
\noindent
In words, after selecting outer sets whose Hausdorff errors decrease
at a prescribed geometric rate, it is sufficient to construct
continuous maps that fix the selected outer sets, approach those outer
sets uniformly, and have summable adjacent uniform displacements.
The resulting subsequence, preceded by the identity map, satisfies
all five clauses of $\mathsf A$.
\begin{proof}
Define the tail displacements
\[
 s_n:=\sum_{j=n}^{\infty}b_j.
\]
The convergence of the series in \eqref{eq:candidate-summability}
implies $s_n\to0$.  Since $a_n\to0$ and $s_n\to0$, one can choose
strictly increasing integers $(j_n)_{n\ge1}$ recursively so that
\[
 j_n\ge n,\qquad
 a_{j_n}\le2^{-n},\qquad
 s_{j_n}\le8\,2^{-n}
 \qquad(n\ge1).
\]
Set
\[
 \widetilde\nu_0:=0,\qquad
 \widetilde R_0:=\id_\D,\qquad
 \widetilde\nu_n:=\nu_{j_n},\qquad
 \widetilde R_n:=T_{j_n}
 \quad(n\ge1).
\]
For $n\ge1$, \eqref{eq:candidate-depths} gives
\[
 \widetilde\nu_n=\nu_{j_n}\ge j_n\ge n,
 \]
and \eqref{eq:candidate-fixing} gives
\[
 \forall c\in O_{\widetilde\nu_n},\qquad
 \widetilde R_n(c)=c.
\]
Moreover, for every $z\in\D$,
\[
 \operatorname{dist}(\widetilde R_n(z),O_{\widetilde\nu_n})
 =\operatorname{dist}(T_{j_n}(z),O_{\nu_{j_n}})
 \le a_{j_n}\le2^{-n}.
\]
Since $O_{\widetilde\nu_n}$ is compact, this distance is attained;
therefore the fourth clause of $\mathsf A$ holds for $n\ge1$.

For $n\ge1$, the triangle inequality gives
\begin{align*}
 \|\widetilde R_{n+1}-\widetilde R_n\|_\infty
 &\le\sum_{j=j_n}^{j_{n+1}-1}
       \|T_{j+1}-T_j\|_\infty\\
 &=\sum_{j=j_n}^{j_{n+1}-1}b_j
 \le s_{j_n}
 \le8\,2^{-n}.
\end{align*}

It remains to verify the clauses at $n=0$.  Since
$p_0(c)=0$, one has $O_0=\overline B(0,2)$.  For every $z\in\D$,
\[
 \operatorname{dist}(z,O_0)\le2\sqrt2-2<1=2^{-0}.
\]
Indeed, this is immediate when $|z|\le2$, and when $|z|>2$ the
point $2z/|z|$ belongs to $O_0$ and is at distance $|z|-2\le
2\sqrt2-2$ from $z$.  Also, $\widetilde R_0=\id_\D$ fixes $O_0$.
Compactness of $O_0$ gives, for each $z\in\D$, a point
$w\in O_0$ with
\[
 |\widetilde R_0(z)-w|
 =\operatorname{dist}(z,O_0)\le2^{-0}.
\]
Finally, both $\widetilde R_0(z)$ and $\widetilde R_1(z)$ belong to
$\D$, so
\[
 \|\widetilde R_1-\widetilde R_0\|_\infty
 \le\operatorname{diam}(\D)=4\sqrt2<8=8\,2^{-0}.
\]
Thus $\mathsf A(\widetilde\nu,\widetilde R_\bullet)$ holds, and
\eqref{eq:ug-existence} gives $\mathsf{UG}$.
\end{proof}

\noindent
Theorem~\ref{thm:outer-gap} and Proposition~\ref{prop:image-gap} are
unconditional.  The convergence $h_N\to0$ is qualitative; no explicit
modulus of the form $N=N(\varepsilon)$ is asserted.  The existence of
maps satisfying \eqref{eq:candidate-fixing}--\eqref{eq:candidate-summability}
is not asserted.  The current implications are
\[
\left[
\begin{gathered}
\text{there exist }\nu\text{ and }T_\bullet\text{ satisfying}\\
\eqref{eq:candidate-depths}\text{--}\eqref{eq:candidate-summability}
\end{gathered}
\right]
\Longrightarrow\mathsf{UG}
\Longrightarrow\LC(\M).
\]

\begin{theorem}[Uniform limit and retraction]\label{thm:uniform}
For $\nu\in\N^\N$ and $R_\bullet\in C(\D,\D)^\N$,
\[
 \mathsf A(\nu,R_\bullet)\Longrightarrow
 \exists R\in C(\D,\D):
 \begin{cases}
  \forall n\in\N,\quad \|R_n-R\|_\infty\le16\,2^{-n},\\
  R(\D)=\M,\\
  \forall c\in\M,\quad R(c)=c.
 \end{cases}
\]
\end{theorem}
\noindent
In words, every pair of sequences satisfying $\mathsf A$ has a
continuous limit map $R:\D\to\D$ whose uniform distance from $R_n$
is at most $16\,2^{-n}$ for every $n$.  The image of $R$ is $\M$,
and $R$ fixes every point of $\M$.
\begin{proof}
Assume $\mathsf A(\nu,R_\bullet)$.  For $m,n\in\N$ with $m>n$, the
triangle inequality and the adjacent-displacement bound give
\begin{align}
 \|R_m-R_n\|_\infty
 &\le\sum_{j=n}^{m-1}\|R_{j+1}-R_j\|_\infty\notag\\
 &\le8\sum_{j=n}^{m-1}2^{-j}\notag\\
 &=16\,2^{-n}\bigl(1-2^{-(m-n)}\bigr)
 \le16\,2^{-n}.\label{eq:cauchy}
\end{align}
For each $z\in\D$, the bound
$|R_m(z)-R_n(z)|\le\|R_m-R_n\|_\infty$ and \eqref{eq:cauchy}
make $(R_n(z))_{n\in\N}$ a Cauchy sequence.  Completeness of $\C$ gives
\[
 R(z):=\lim_{n\to\infty}R_n(z).
\]
Since $\D$ is closed and each $R_n(z)$ belongs to $\D$, the limit
$R(z)$ belongs to $\D$.  For fixed $n$ and $z$, the pointwise bound
$|R_m(z)-R_n(z)|\le16\,2^{-n}$ holds for every $m>n$.
Taking $m\to\infty$ and using continuity of the modulus yields
\begin{equation}\label{eq:tail}
 \forall n\in\N\ \forall z\in\D,\qquad
 |R_n(z)-R(z)|\le16\,2^{-n}.
\end{equation}

For continuity, fix $z_0\in\D$ and $\varepsilon>0$.  Choose $n\in\N$
such that $32\,2^{-n}<\varepsilon/2$.  Continuity of $R_n$ gives
$\delta>0$ such that, for every $z\in\D$ with $|z-z_0|<\delta$,
$|R_n(z)-R_n(z_0)|<\varepsilon/2$.  For these $z$,
\begin{align*}
 |R(z)-R(z_0)|
 &\le |R(z)-R_n(z)|+|R_n(z)-R_n(z_0)|+|R_n(z_0)-R(z_0)|\\
 &\le32\,2^{-n}+|R_n(z)-R_n(z_0)|<\varepsilon.
\end{align*}
Thus $R\in C(\D,\D)$.  Taking the supremum over $z$ in
\eqref{eq:tail} gives $\|R_n-R\|_\infty\le16\,2^{-n}$; hence the
convergence is uniform.

Fix $z\in\D$ and $k\in\N$.  For every $n\ge k$, choose
$w_n\in O_{\nu_n}$ with $|R_n(z)-w_n|\le2^{-n}$.
Since $\nu_n\ge n\ge k$, Lemma~\ref{lem:outer} gives
$O_{\nu_n}\subseteq O_k$.  Moreover,
\[
 |R(z)-w_n|
 \le|R(z)-R_n(z)|+|R_n(z)-w_n|
 \le17\,2^{-n}\longrightarrow0.
\]
The set $O_k$ is closed, so $R(z)\in O_k$.  As $k$ was arbitrary,
Lemma~\ref{lem:outer} implies
$R(z)\in\bigcap_{k\in\N}O_k=\M$.

For $c\in\M$, the same lemma gives $c\in O_{\nu_n}$ for every $n$.
The fixing clause in \eqref{eq:hypotheses} therefore gives
$R_n(c)=c$ for every $n$, and taking limits gives $R(c)=c$.
Consequently $R(\D)\subseteq\M$ and, because $\M\subseteq\D$,
$\M\subseteq R(\D)$.  Thus $R(\D)=\M$.
\end{proof}

\section{Main current result}

\begin{theorem}[Conditional Mandelbrot local connectivity]\label{thm:mlc}
\[
 \boxed{\mathsf{UG}\Longrightarrow\LC(\M).}
\]
\end{theorem}
\noindent
In words, if there exist sequences $\nu$ and $R_\bullet$ satisfying
$\mathsf A(\nu,R_\bullet)$, then for every $c\in\M$ and every
$\eta>0$ there exist $\delta>0$ and a connected subset $V$ of $\M$
with
$\M\cap B(c,\delta)\subseteq V\subseteq\M\cap B(c,\eta)$.
By Lemma~\ref{lem:lc-criterion}, this conclusion states that $\M$
is locally connected.
\begin{proof}
Assume $\mathsf{UG}$ and choose sequences satisfying
$\mathsf A(\nu,R_\bullet)$ by \eqref{eq:ug-existence}.
Theorem~\ref{thm:uniform} gives a continuous map $R:\D\to\D$
such that $R(\D)=\M$ and $R(c)=c$ for every $c\in\M$.

Fix $c\in\M$ and $\eta>0$.  Since $R(c)=c$, continuity at $c$ gives
$\delta>0$ satisfying
\[
 \forall z\in\D,\qquad
 |z-c|<\delta\Longrightarrow |R(z)-c|<\eta.
\]
Define
\[
 W:=\D\cap B(c,\delta),\qquad V:=R(W).
\]
The square $\D$ and the disk $B(c,\delta)$ are convex, so $W$ is
convex.  It contains $c$ and is therefore nonempty and connected.
Its continuous image $V$ is connected, and $V\subseteq\M$ because
$R(\D)=\M$.

If $y\in\M\cap B(c,\delta)$, then $y\in W$ and $R(y)=y$, so
$y\in V$.  If $y\in V$, choose $z\in W$ with $y=R(z)$.
Then $y\in\M$ and $|y-c|=|R(z)-c|<\eta$.  Hence
\[
 \M\cap B(c,\delta)\subseteq V\subseteq\M\cap B(c,\eta).
\]
The choices were made for arbitrary $c\in\M$ and $\eta>0$, proving
$\LC(\M)$.  Lemma~\ref{lem:lc-criterion} identifies this with local
connectedness in the Euclidean subspace topology.
\end{proof}

\end{document}
